\clearpage
\fancyhead[L]{\small MERCURIUS IN SOLE VISUS.}
\hypertarget{native-653-lower}{}
\begin{center}\small Göttingische gelehrte Anzeigen. No. 169. p. 1688. 21 October 1833.\\[3pt]\rule{.72\linewidth}{.3pt}\end{center}

\emph{Mercurius in Sole visus, of Overgang van Mercurius over de Zon; den 5 Mei \textup{1832} te Utrecht waargenomen door} G. \textsc{Moll}. Amsterdam 1833. Published by C. G. Sulpke. 43 pages. 4to.

In this short work we find the observations of the most recent transit of Mercury in Utrecht at the observatory, and at the home of Mr. \textsc{Van Beek}, as well as in Leiden, Amsterdam and Nimwegen; further, a compilation of all observations of earlier transits made in Holland or by Dutch astronomers; finally, the observations made in Utrecht and Leyden of the occultation of Saturn by the Moon on 8 May 1832. A more detailed review would be superfluous here, since the substance of its contents has already been reported in issue 229 of Astronomische Nachrichten, on the basis of letters from Prof. \textsc{Moll} to State Councillor \textsc{Schumacher}.
